% paper_appendix.tex -- the experimental-details appendix, to paste into the paper.
% Requires: amsmath, amssymb, booktabs, xcolor, listings, tikz
% (\usetikzlibrary{positioning,arrows.meta,calc}), algorithm. Uses the paper's \pranav and
% \jeffrey comment macros.
%
% Notation is that of Secs. 3-5 and Apps. A-B: bold vectors; seed z ~ gamma = N(0, I_d); model time
% tau (tau = 1 seed, tau = 0 data); sampling fields F_star (exact) and F_theta (learned) with
% endpoint maps G_star, G_theta = Phi_{1->0}; mixture weights w_i, centres mu_i, component
% covariance sigma_0^2 I_d; mode core C_k^bullet = B(mu_k, r_k) with r_k = R99 (Sec. 4); pulled-back
% core K_k^F = G_F^{-1}(C_k^bullet); eta_t = sqrt((1 - abar_t)/abar_t) the added noise scale of
% Sec. B. Symbols the paper reserves (pi_i, sigma(t), u, h_i^F, rho_i, D_sigma, S_ij, B(c,a), E_c,
% theta, kappa_c, o(.), a(u), x_a, z_a) are not reused with another meaning.
% Algorithm 1 (pullback) is the listing of Sec. 4 (tables/algos/pullback.tex); it is referred to as
% Algorithm 1 and not input again here. Listing identifiers follow Algorithm 1 (mu, R99, n, m, w,
% F_star, T, Z, y); every identifier a listing uses is named in its header, and the constants are
% valued in Tables app-grid, app-scorenet, app-classifiers and app-mnist.
% Every statement was checked against the code, configs, checkpoints and logs of run abc123:
%   Gaussian mixtures -- failsafe_iclr2 (code/, conf/, output/abc123/, checkpoints/), d = 2..512
%   (d = 512 evaluated on Euler, merged locally on 2026-09-25);
%   image latents -- ~/Desktop/iclr/experiment-2, run abc123 (training and evaluation with the code
%   of commit 2be1abf, results in 24c85fc; decoded-sample grids with commit 746f4e1).

\definecolor{appcodecomment}{RGB}{38,139,150}
\definecolor{appcodekeyword}{RGB}{0,70,170}
\definecolor{appcodecall}{RGB}{190,30,130}
\lstdefinestyle{appcode}{
  language=Python, basicstyle=\ttfamily\footnotesize, columns=fullflexible, keepspaces=true,
  commentstyle=\color{appcodecomment}, keywordstyle=\color{appcodekeyword},
  emphstyle=\color{appcodecall}, showstringspaces=false, aboveskip=2pt, belowskip=0pt,
  morekeywords={lambda},
  emph={range, len, max, min, sum, mean, log, exp, sqrt, ceil, norm, cdist, topk, top2, argmax,
        argmin, where, one_hot, softmax, log_softmax, entropy, concat, triu_indices, quantiles,
        argmax_over, randn, unit_vectors, pullback, transport, outcome, train, train_field,
        draw_centres, draw_weights, sample_gmm, hallucination_rate, metrics, hall_f1, AdamW,
        Linear, cross_entropy, G_theta, G_star, G_star_hat, train_autoencoder, R99_per_cluster,
        invert, Classifier, fit, predict, kmeans_mixture, Enc},
}


\section{Experimental Details}
\label{app:experimental_details}

This appendix specifies the experiments of Sec.~5, in which pullback predicts, before sampling,
the outcome of every seed $\boldsymbol{z}$: the mode core $\mathcal{C}_k^\bullet$ that its sample
reaches, or a hallucination. We describe the Gaussian-mixture targets and their mode cores
(Sec.~\ref{app:details_gmm}), the pullback anchors and the evaluation protocol
(Sec.~\ref{app:anchor-class}), the learned sampling fields whose outcomes are predicted
(Sec.~\ref{app:details-diffusion}), and the anchor classifiers (Sec.~\ref{app:algo-details}).
Sec.~\ref{app:mnist} extends the setup to the latent space of an image autoencoder.


\subsection{Gaussian Mixture Setup}
\label{app:details_gmm}

\paragraph{Target distribution.}
The data distribution is a mixture of $K$ isotropic Gaussians in $\mathbb{R}^d$,
\begin{equation}
    p_{\text{data}}(\boldsymbol{x}) = \sum_{i=1}^{K} w_i \, \mathcal{N}(\boldsymbol{x}; \boldsymbol{\mu}_i, \sigma_0^2 \boldsymbol{I}_d),
    \qquad \sigma_0 = 0.1 .
\end{equation}
Let $R_{99} = \sigma_0 \sqrt{\chi^2_d(0.99)}$, where $\chi^2_d(0.99)$ is the $0.99$-quantile of the
chi-square distribution with $d$ degrees of freedom, denote the radius of the ball
$B(\boldsymbol{\mu}_i, R_{99})$ that holds $99\%$ of the mass of component $i$; it grows from
$R_{99} = 0.30$ at $d = 2$ to $1.77$ at $d = 256$ and $2.43$ at $d = 512$. The centres are placed
one at a time: each proposal $2\sqrt{d/2}\,\boldsymbol{g}/\|\boldsymbol{g}\|_2$,
$\boldsymbol{g} \sim \mathcal{N}(\boldsymbol{0}, \boldsymbol{I}_d)$, is uniform on the sphere of
radius $2\sqrt{d/2}$ and is accepted if it lies at least $3R_{99}$ from every centre already
placed, and redrawn otherwise, so that every pair satisfies
$d_{ij} = \|\boldsymbol{\mu}_i - \boldsymbol{\mu}_j\|_2 \ge 3R_{99}$. For $d = 2$ at most 13 centres
fit on the circle of radius $2$ at this separation (16 points on it are at most
$4\sin(\pi/16) = 0.78 < 3R_{99} = 0.91$ apart), so the pair $(K, d) = (16, 2)$ is not evaluated.
We consider an \emph{unweighted} mixture, $w_i = 1/K$, and a \emph{weighted} one, for which $K$
weights are drawn once from $U[0.3, 0.8]$ with a fixed seed and normalised; this profile depends
only on $K$, and each run assigns it to the centres in its own random order. The two variants are
separate experiments with separately trained models.

\paragraph{Mode cores and outcomes.}
Following Sec.~4, the critical radius of a converged sample is $R_{99}$: the mode core of
component $k$ is $\mathcal{C}_k^\bullet = B(\boldsymbol{\mu}_k, r_k)$ with $r_k = R_{99}$ for every
$k$, and the separation $3R_{99} > 2R_{99}$ makes the cores disjoint. For this mixture the exact
sampling field $\boldsymbol{F}_\star$ has a closed form, because the noised mixture is again a
Gaussian mixture; the weights $w_i$ enter through its posterior over components. For DDIM and the
solvers that reuse its network, $\boldsymbol{F}_\star$ is the probability-flow field of the
variance-preserving process of Sec.~B with the exact score of the noised mixture; for flow
matching, it is the marginal velocity of the optimal-transport conditional path of
Sec.~\ref{app:details-diffusion}. A trained network gives the learned field
$\boldsymbol{F}_\theta$. Every sampler we use is deterministic, so a seed
$\boldsymbol{z} \sim \gamma$ has a single sample
$\boldsymbol{x} = \mathcal{G}_{\boldsymbol{F}}(\boldsymbol{z})$, where
$\mathcal{G}_{\boldsymbol{F}} = \Phi^{\boldsymbol{F}}_{1 \to 0}$ is the endpoint map of Eq.~(1) for
$\boldsymbol{F} \in \{\boldsymbol{F}_\star, \boldsymbol{F}_\theta\}$, and we label the seed by its
\emph{outcome}
\begin{equation}
    y(\boldsymbol{z}) =
    \begin{cases}
        k,  & \mathcal{G}_{\boldsymbol{F}}(\boldsymbol{z}) \in \mathcal{C}_k^\bullet, \\
        -1, & \mathcal{G}_{\boldsymbol{F}}(\boldsymbol{z}) \notin \bigcup_{k} \mathcal{C}_k^\bullet
              \quad \text{(hallucination)}.
    \end{cases}
\end{equation}
Equivalently, $y(\boldsymbol{z}) = k$ exactly when $\boldsymbol{z}$ lies in the pulled-back core
$\mathcal{K}_k^{\boldsymbol{F}} = \mathcal{G}_{\boldsymbol{F}}^{-1}(\mathcal{C}_k^\bullet)$, and the
hallucinating seeds form the hallucination band outside every $\mathcal{K}_k^{\boldsymbol{F}}$.
Predicting the outcome from the seed alone is therefore a $(K+1)$-class classification problem in
seed space whose decision boundaries are the pulled-back boundaries
$\partial\mathcal{K}_k^{\boldsymbol{F}}$.


\subsection{Anchor Classification}
\label{app:anchor-class}

\paragraph{Pullback anchors.}
The pulled-back cores have no closed form, so we describe them by \emph{anchors}: seeds of known
outcome, obtained by pullback of labelled points in data space (Algorithm~1). For a budget of $b$
anchors per mode, we draw $n = b$ points uniformly in each mode core $\mathcal{C}_k^\bullet$
(label $k$) and $m = b/2$ points in the data-space band
$R_{99} < \|\boldsymbol{x} - \boldsymbol{\mu}_k\|_2 \le R_{99} + w$ around it (label $-1$), with
band width $w = 2\sigma_0$. Each point is
$\boldsymbol{x} = \boldsymbol{\mu}_k + r\, \boldsymbol{g}/\|\boldsymbol{g}\|_2$, where
$\boldsymbol{g} \sim \mathcal{N}(\boldsymbol{0}, \boldsymbol{I}_d)$ makes the direction uniform on
the unit sphere and the radius $r$ is drawn with $U \sim U[0, 1]$ as
\begin{equation}
    r = R_{99}\, U^{1/d} \quad \text{(core)}, \qquad
    r = R_{99} + w\, U \quad \text{(band)}.
\end{equation}
The exponent $1/d$ makes the core points uniform in volume. The band points are uniform in radius,
which keeps them spread across the band in high dimension, where a volume-uniform draw would
concentrate them at its outer edge. They supply hallucination labels directly outside the core
boundaries $\partial\mathcal{C}_k^\bullet$, whose pullbacks
$\partial\mathcal{K}_k^{\boldsymbol{F}_\star}$ the classifiers must locate. The band has a fixed
width while $R_{99}$ grows with $d$, so it narrows relative to the core, from $w/R_{99} = 0.66$ at
$d = 2$ to $0.08$ at $d = 512$. The offsets of the anchors from the centres are drawn with the same
seed in every run.

The labelled points are pulled back by integrating $\boldsymbol{F}_\star$ from $\tau = 0$ to
$\tau = 1$ on a grid of $T$ time points, i.e.\ in $T - 1$ steps (Figure~\ref{fig:pipeline}, top).
For flow matching, the exact velocity is integrated with Heun's method on the uniform grid, so $T$
changes only the discretisation. For DDIM, each step is a DDIM update whose noise prediction
averages the exact $\boldsymbol{\epsilon}$ at both ends of the step, a Heun-type correction that
makes the pullback second order (the learned DDIM sampler uses the plain first-order update). The
DDIM pullback uses the linear schedule $\beta \in [10^{-4}, 0.02]$ discretised into $T$ noise
levels, $\bar\alpha_t = \prod_{s \le t}(1 - \beta_s)$ with $\beta_1, \dots, \beta_T$ equally
spaced, so its terminal $\bar\alpha_T$ equals that of the learned sampler (the same schedule with
500 levels) only for $T = 500$ ($\bar\alpha_T = 8.0\times10^{-2}$, $6.4\times10^{-3}$ and
$5.1\times10^{-4}$ for $T = 250, 500, 750$). Pullback never queries the learned model.

\paragraph{Reference outcomes.}
The test set consists of $N K$ seeds $\boldsymbol{z} \sim \gamma$, $N = 50{,}000$, labelled by
their outcome under the learned endpoint map $\mathcal{G}_\theta$ (Figure~\ref{fig:pipeline},
bottom). The anchors come from $\boldsymbol{F}_\star$ and the reference outcomes from
$\boldsymbol{F}_\theta$, so the evaluation measures how well the pulled-back cores of the exact
field predict the trained sampler. The choice of $T$ affects only the exact-field side of the
evaluation: the anchors and the exact-field labels of the w-$k$NN calibration seeds
(Sec.~\ref{app:algo-details}).

\begin{figure}[htbp]
\centering
\begin{tikzpicture}[
    font=\footnotesize, >=Stealth, node distance=4mm and 4mm,
    blk/.style={draw, rounded corners=2pt, align=center, text width=24mm, inner sep=3pt, minimum height=15mm},
    top/.style={blk, fill=blue!8}, bot/.style={blk, fill=orange!12},
    res/.style={blk, fill=green!10, text width=20mm}]
  \node[top] (a1) {Sample $n$ points per mode core $\mathcal{C}_k^\bullet$ (label $k$), $m$ in the band (label $-1$)};
  \node[top, right=of a1] (a2) {Pullback with the \emph{exact} field $\boldsymbol{F}_\star$ on a $T$-point grid};
  \node[top, right=of a2] (a3) {Labelled anchors in seed space};
  \node[top, right=of a3] (a4) {Fit each classifier on the anchors (w-$k$NN: its own)};
  \node[bot, below=8mm of a1] (g1) {Draw $N K$ test seeds $\boldsymbol{z} \sim \mathcal{N}(\boldsymbol{0}, \boldsymbol{I}_d)$};
  \node[bot, right=of g1] (g2) {Generate $\boldsymbol{x} = \mathcal{G}_\theta(\boldsymbol{z})$ on a 500-point grid};
  \node[bot, right=of g2] (g3) {Outcome: $k$ if $\boldsymbol{x} \in \mathcal{C}_k^\bullet$, else $-1$};
  \node[bot, right=of g3] (g4) {Reference outcomes of the test seeds};
  \node[res, right=10mm of $(a4.east)!0.5!(g4.east)$] (res) {Compare:\\accuracy,\\mode F1,\\hallucination F1};
  \draw[->] (a1) -- (a2);
  \draw[->] (a2) -- (a3);
  \draw[->] (a3) -- (a4);
  \draw[->] (g1) -- (g2);
  \draw[->] (g2) -- (g3);
  \draw[->] (g3) -- (g4);
  \draw[->] (a4.east) -| node[pos=0.25, above] {predicted} (res.north);
  \draw[->] (g4.east) -| node[pos=0.25, below] {reference} (res.south);
\end{tikzpicture}
\caption{\textbf{Evaluation protocol.} Top: points in and around each mode core are pulled back to
seed space with the exact field $\boldsymbol{F}_\star$, and the classifiers are fit on the resulting
anchors (w-$k$NN on its own pulled-back, mode-labelled anchors, with its threshold calibrated on
seeds labelled by $\mathcal{G}_\star$; Sec.~\ref{app:algo-details}). Bottom: test seeds are labelled
by their outcome under the learned endpoint map $\mathcal{G}_\theta$. The classifiers are evaluated
on the same test seeds.}
\label{fig:pipeline}
\end{figure}

\paragraph{Metrics.}
We report the accuracy over the $K + 1$ outcome classes, the F1 score macro-averaged over the mode
classes, and the F1 score of the hallucination class. Hallucinations are rare: averaged over the
$(K, d)$ grid and the runs, the learned samplers send $1.0\%$ (DDIM) and $0.8\%$ (flow matching) of
the test seeds outside every mode core, close to the $1\%$ of its mass that $p_{\text{data}}$ itself
places outside the cores, and $1.6$--$1.8\%$ under Heun, RK45 and DPM-Solver++(2M), which reuse the
DDIM network (up to $3.7\%$ in a single run at $d = 512$). Accuracy is therefore dominated by seeds
inside the pulled-back cores, and hallucination F1 is the relevant measure of how well the
hallucination band is recovered. All results are means over three runs (seeds 0, 100 and 200).
Each run draws its own centres, weight assignment, training set, network initialisation and test
seeds, while the anchor offsets and the classifier initialisations use the same seeds in every run.
The test seeds and the centres depend only on $(K, d)$ and the run, so all samplers and both
mixture variants are evaluated on the same seeds and centres, and comparisons between them are
paired.

\paragraph{Protocol.}
An experiment is specified by the sampler, the mixture variant, $K$, $d$, the number of pullback
steps $T$ and the anchor budget $b$. Algorithm~\ref{alg:app-protocol} gives the procedure for one
configuration, and Table~\ref{tab:app-grid} lists the values we sweep; every combination is
evaluated except $(K, d) = (16, 2)$. The constants named in the algorithms are valued in
Tables~\ref{tab:app-grid}, \ref{tab:app-scorenet} and~\ref{tab:app-classifiers}. The learned field
and the reference outcomes depend on neither $T$ nor $b$, so they are computed once per run and
shared by all values of $T$ and $b$; Heun, RK45 and DPM-Solver++(2M) moreover reuse the DDIM
network. Each result table fixes the sampler, the mixture variant and $b$ and reports the overall
accuracy (\%, mean over the three runs): each row is a pair $(K, d)$, each classifier has one
sub-column per value of $T$, and the best classifier at each $T$ is in bold. The heatmaps, one per
sampler, variant, $T$ and $b$, show the three metrics over the $(K, d)$ grid, one panel per
classifier and metric. Tables~\ref{tab:pullback-T-sweep} and~\ref{tab:pullback-T-sweep-flow}
summarise the unweighted DDIM and flow-matching results at the largest budget. For flow matching
$T$ only refines the discretisation, and the accuracies barely change with it. For DDIM, where only
$T = 500$ reproduces the schedule of the learned sampler, the polar classifier averages $94.3\%$ and
$95.1\%$ over $d \ge 128$ at $T = 500$ and $750$, but $49.6\%$ at $T = 250$.

% Auto-generated by experiment_results/build_table_s1.py -- needs booktabs

\begin{table}[p]
\centering
\scriptsize
\setlength{\tabcolsep}{3pt}
\caption{\textbf{Outcome prediction across pullback grids, DDIM.} Overall accuracy (\%, mean over 3 runs) of each classifier fit on $b = 20{,}000$ anchors per mode, for $T = 250, 500, 750$ pullback grid points (sub-columns), on the unweighted mixture; hall.\ is the hallucination rate (\%) of the learned sampler on the test seeds, which does not depend on $T$. Best classifier per $(K, d, T)$ in bold; $(K, d) = (16, 2)$ is not evaluated (Sec.~\ref{app:details_gmm}).}
\label{tab:pullback-T-sweep}
\begin{tabular}{rrrcccccccccccc}
\toprule
 & & & \multicolumn{3}{c}{$k$NN} & \multicolumn{3}{c}{w-$k$NN} & \multicolumn{3}{c}{quadratic} & \multicolumn{3}{c}{polar} \\
\cmidrule(lr){4-6} \cmidrule(lr){7-9} \cmidrule(lr){10-12} \cmidrule(lr){13-15}
$K$ & $d$ & hall. & 250 & 500 & 750 & 250 & 500 & 750 & 250 & 500 & 750 & 250 & 500 & 750 \\
\midrule
2 & 2 & 1.0 & \textbf{99.0} & \textbf{99.5} & \textbf{99.5} & 98.6 & 98.5 & 98.4 & 98.1 & 98.8 & 98.9 & \textbf{99.0} & \textbf{99.5} & \textbf{99.5} \\
2 & 4 & 0.9 & 95.2 & 95.2 & 95.2 & 90.9 & 91.1 & 91.4 & 98.2 & 98.9 & 99.0 & \textbf{98.5} & \textbf{99.5} & \textbf{99.4} \\
2 & 8 & 1.1 & 86.5 & 85.6 & 85.5 & 83.5 & 90.4 & 84.6 & \textbf{97.6} & 98.9 & 98.8 & \textbf{97.6} & \textbf{99.4} & \textbf{99.4} \\
2 & 16 & 1.1 & 81.5 & 80.5 & 80.4 & 80.1 & 83.5 & 85.6 & 96.3 & 98.7 & 98.6 & \textbf{96.9} & \textbf{99.3} & \textbf{99.2} \\
2 & 32 & 1.0 & 81.5 & 77.0 & 76.7 & 76.7 & 81.8 & 80.1 & 89.0 & 98.3 & 98.3 & \textbf{94.5} & \textbf{99.2} & \textbf{99.2} \\
2 & 64 & 1.0 & 82.5 & 73.9 & 72.8 & 61.4 & 84.5 & 80.3 & 51.0 & 96.5 & 97.0 & \textbf{87.2} & \textbf{98.9} & \textbf{99.0} \\
2 & 128 & 0.9 & \textbf{84.9} & 71.7 & 68.7 & 61.9 & 73.7 & 72.8 & 9.3 & 88.3 & 91.3 & 70.6 & \textbf{97.6} & \textbf{97.8} \\
2 & 256 & 0.9 & \textbf{86.0} & 70.0 & 64.8 & 62.7 & 60.8 & 54.6 & 21.7 & 41.9 & 47.2 & 46.7 & \textbf{96.7} & \textbf{97.5} \\
2 & 512 & 0.9 & \textbf{87.3} & 70.0 & 61.7 & 64.3 & 63.3 & 47.3 & 18.0 & 35.2 & 39.3 & 26.3 & \textbf{94.5} & \textbf{96.7} \\
\midrule
4 & 2 & 1.1 & \textbf{92.4} & \textbf{98.9} & \textbf{97.3} & 91.9 & 98.1 & 96.7 & 90.8 & 93.2 & 92.7 & 92.1 & 98.2 & 96.6 \\
4 & 4 & 1.1 & 89.7 & 90.9 & 90.6 & 87.4 & 89.1 & 90.0 & 91.6 & 94.0 & 93.1 & \textbf{92.5} & \textbf{97.9} & \textbf{96.7} \\
4 & 8 & 1.1 & 75.9 & 74.8 & 74.4 & 85.1 & 86.5 & 86.6 & 94.0 & 95.4 & 94.9 & \textbf{94.7} & \textbf{98.3} & \textbf{97.6} \\
4 & 16 & 1.0 & 67.2 & 65.2 & 65.2 & 83.2 & 82.6 & 80.7 & 95.5 & 96.8 & 96.6 & \textbf{96.6} & \textbf{98.1} & \textbf{97.7} \\
4 & 32 & 1.0 & 65.6 & 60.2 & 59.9 & 83.0 & 80.8 & 81.9 & 92.2 & 96.8 & 96.8 & \textbf{95.3} & \textbf{98.3} & \textbf{97.8} \\
4 & 64 & 1.0 & 67.3 & 55.1 & 54.0 & 77.7 & 76.8 & 71.9 & 68.9 & 94.5 & 94.8 & \textbf{90.3} & \textbf{97.6} & \textbf{97.3} \\
4 & 128 & 0.8 & \textbf{70.9} & 50.8 & 48.1 & 56.3 & 71.9 & 67.3 & 17.3 & 86.2 & 88.1 & 69.4 & \textbf{97.0} & \textbf{97.1} \\
4 & 256 & 0.7 & \textbf{74.1} & 47.9 & 42.7 & 58.8 & 66.6 & 63.1 & 17.9 & 39.6 & 48.3 & 46.6 & \textbf{95.6} & \textbf{96.5} \\
4 & 512 & 0.8 & \textbf{75.7} & 46.8 & 38.3 & 59.6 & 56.0 & 46.4 & 12.1 & 23.1 & 27.7 & 22.1 & \textbf{91.6} & \textbf{95.3} \\
\midrule
8 & 2 & 1.1 & \textbf{96.7} & \textbf{98.1} & \textbf{97.9} & 96.5 & 97.5 & 97.3 & 83.2 & 84.2 & 84.3 & 96.3 & 97.0 & 96.7 \\
8 & 4 & 1.2 & 86.8 & 88.4 & 87.7 & 86.4 & 89.3 & 87.4 & 88.8 & 90.5 & 89.6 & \textbf{90.4} & \textbf{96.3} & \textbf{95.1} \\
8 & 8 & 1.2 & 68.1 & 67.4 & 67.2 & 81.1 & 82.8 & 82.6 & 92.2 & 93.2 & 92.8 & \textbf{94.4} & \textbf{97.2} & \textbf{96.5} \\
8 & 16 & 1.1 & 55.4 & 53.6 & 53.5 & 76.3 & 76.8 & 77.5 & 93.0 & 94.0 & 93.5 & \textbf{95.2} & \textbf{97.2} & \textbf{96.2} \\
8 & 32 & 1.1 & 51.9 & 46.6 & 46.4 & 81.2 & 76.8 & 75.7 & 92.3 & 94.8 & 94.5 & \textbf{94.4} & \textbf{97.0} & \textbf{96.2} \\
8 & 64 & 1.1 & 53.2 & 40.6 & 39.9 & 82.8 & 75.1 & 72.1 & 80.1 & 92.6 & 92.6 & \textbf{91.5} & \textbf{96.2} & \textbf{95.6} \\
8 & 128 & 0.8 & 57.9 & 35.5 & 33.6 & 56.8 & 62.7 & 65.3 & 29.1 & 84.5 & 85.3 & \textbf{75.3} & \textbf{95.9} & \textbf{95.8} \\
8 & 256 & 0.7 & \textbf{61.8} & 31.6 & 28.0 & 53.2 & 46.0 & 48.8 & 14.6 & 40.4 & 49.9 & 50.3 & \textbf{95.2} & \textbf{95.5} \\
8 & 512 & 0.6 & \textbf{64.6} & 29.8 & 23.6 & 51.8 & 57.8 & 39.1 & 7.4 & 19.5 & 27.0 & 25.2 & \textbf{91.6} & \textbf{93.8} \\
\midrule
16 & 4 & 1.3 & 86.5 & 87.5 & 87.2 & 87.8 & 88.7 & 88.9 & 88.4 & 89.5 & 89.2 & \textbf{92.2} & \textbf{95.1} & \textbf{94.5} \\
16 & 8 & 1.3 & 61.7 & 61.1 & 60.9 & 78.7 & 77.3 & 79.6 & 90.4 & 91.3 & 90.5 & \textbf{91.9} & \textbf{95.4} & \textbf{94.6} \\
16 & 16 & 1.2 & 45.9 & 44.9 & 44.9 & 73.6 & 72.1 & 72.0 & 92.7 & 93.8 & 93.3 & \textbf{94.4} & \textbf{96.0} & \textbf{95.5} \\
16 & 32 & 1.1 & 40.2 & 36.1 & 36.1 & 78.1 & 72.3 & 71.4 & 92.4 & 94.3 & 94.1 & \textbf{93.9} & \textbf{95.6} & \textbf{95.2} \\
16 & 64 & 1.2 & 40.6 & 29.9 & 29.5 & 80.1 & 69.9 & 69.5 & 86.4 & 92.5 & 92.5 & \textbf{90.1} & \textbf{95.2} & \textbf{94.8} \\
16 & 128 & 0.7 & 45.2 & 24.4 & 23.3 & 59.0 & 57.7 & 53.1 & 45.1 & 84.9 & 85.2 & \textbf{79.0} & \textbf{94.5} & \textbf{94.0} \\
16 & 256 & 0.6 & 50.3 & 20.4 & 18.2 & 50.9 & 44.0 & 42.8 & 12.3 & 41.9 & 50.5 & \textbf{53.6} & \textbf{92.8} & \textbf{92.6} \\
16 & 512 & 1.4 & \textbf{53.5} & 18.2 & 14.4 & 52.3 & 45.4 & 34.0 & 5.2 & 18.9 & 29.5 & 29.6 & \textbf{88.3} & \textbf{89.1} \\
\bottomrule
\end{tabular}
\end{table}

\begin{table}[p]
\centering
\scriptsize
\setlength{\tabcolsep}{3pt}
\caption{\textbf{Outcome prediction across pullback grids, flow matching.} Overall accuracy (\%, mean over 3 runs) of each classifier fit on $b = 20{,}000$ anchors per mode, for $T = 250, 500, 750$ pullback grid points (sub-columns), on the unweighted mixture; hall.\ is the hallucination rate (\%) of the learned sampler on the test seeds, which does not depend on $T$. Best classifier per $(K, d, T)$ in bold; $(K, d) = (16, 2)$ is not evaluated (Sec.~\ref{app:details_gmm}).}
\label{tab:pullback-T-sweep-flow}
\begin{tabular}{rrrcccccccccccc}
\toprule
 & & & \multicolumn{3}{c}{$k$NN} & \multicolumn{3}{c}{w-$k$NN} & \multicolumn{3}{c}{quadratic} & \multicolumn{3}{c}{polar} \\
\cmidrule(lr){4-6} \cmidrule(lr){7-9} \cmidrule(lr){10-12} \cmidrule(lr){13-15}
$K$ & $d$ & hall. & 250 & 500 & 750 & 250 & 500 & 750 & 250 & 500 & 750 & 250 & 500 & 750 \\
\midrule
2 & 2 & 1.0 & 99.6 & 99.6 & 99.6 & 98.6 & 98.6 & 98.6 & 98.9 & 98.9 & 98.9 & \textbf{99.7} & \textbf{99.7} & \textbf{99.7} \\
2 & 4 & 0.9 & 95.2 & 95.2 & 95.2 & 93.7 & 93.8 & 93.8 & 99.0 & 99.0 & 99.0 & \textbf{99.5} & \textbf{99.5} & \textbf{99.5} \\
2 & 8 & 1.0 & 85.5 & 85.5 & 85.5 & 86.5 & 86.5 & 86.5 & 98.9 & 99.0 & 98.9 & \textbf{99.5} & \textbf{99.5} & \textbf{99.5} \\
2 & 16 & 0.9 & 80.2 & 80.2 & 80.2 & 81.5 & 82.5 & 82.5 & 98.8 & 98.8 & 98.8 & \textbf{99.4} & \textbf{99.4} & \textbf{99.4} \\
2 & 32 & 0.8 & 76.7 & 76.6 & 76.6 & 76.8 & 78.0 & 78.0 & 98.5 & 98.4 & 98.4 & \textbf{99.4} & \textbf{99.4} & \textbf{99.4} \\
2 & 64 & 0.8 & 72.5 & 72.7 & 72.7 & 81.7 & 80.8 & 80.8 & 97.1 & 97.1 & 97.1 & \textbf{99.1} & \textbf{99.2} & \textbf{99.2} \\
2 & 128 & 0.6 & 68.6 & 68.5 & 68.5 & 72.1 & 63.5 & 63.5 & 91.8 & 91.7 & 91.7 & \textbf{97.9} & \textbf{98.0} & \textbf{98.0} \\
2 & 256 & 0.4 & 64.4 & 64.4 & 64.4 & 58.8 & 58.8 & 58.8 & 48.4 & 48.4 & 48.4 & \textbf{97.6} & \textbf{97.7} & \textbf{97.6} \\
2 & 512 & 0.3 & 60.8 & 60.8 & 60.9 & 58.1 & 57.9 & 65.0 & 39.8 & 40.4 & 40.0 & \textbf{97.1} & \textbf{97.2} & \textbf{97.1} \\
\midrule
4 & 2 & 1.0 & \textbf{98.9} & \textbf{98.9} & \textbf{98.9} & 98.2 & 98.3 & 98.3 & 93.4 & 93.4 & 93.4 & 98.2 & 98.2 & 98.2 \\
4 & 4 & 1.0 & 90.8 & 90.8 & 90.8 & 89.6 & 89.6 & 89.6 & 94.2 & 94.2 & 94.2 & \textbf{98.1} & \textbf{98.1} & \textbf{98.1} \\
4 & 8 & 1.1 & 74.5 & 74.5 & 74.5 & 87.4 & 87.4 & 87.4 & 95.5 & 95.5 & 95.5 & \textbf{98.4} & \textbf{98.5} & \textbf{98.4} \\
4 & 16 & 0.9 & 65.1 & 65.1 & 65.1 & 82.3 & 82.3 & 82.3 & 97.0 & 97.0 & 97.0 & \textbf{98.1} & \textbf{98.1} & \textbf{98.1} \\
4 & 32 & 1.0 & 59.8 & 59.7 & 59.7 & 80.1 & 81.7 & 81.7 & 96.7 & 96.7 & 96.7 & \textbf{98.4} & \textbf{98.5} & \textbf{98.5} \\
4 & 64 & 0.9 & 53.9 & 53.8 & 53.8 & 75.9 & 78.6 & 78.6 & 95.0 & 95.0 & 95.0 & \textbf{97.8} & \textbf{97.9} & \textbf{97.9} \\
4 & 128 & 0.6 & 47.9 & 47.9 & 47.9 & 65.6 & 61.7 & 61.7 & 88.4 & 88.4 & 88.4 & \textbf{97.7} & \textbf{97.5} & \textbf{97.5} \\
4 & 256 & 0.4 & 42.2 & 42.2 & 42.2 & 60.4 & 60.4 & 60.4 & 49.7 & 49.7 & 49.7 & \textbf{97.0} & \textbf{97.0} & \textbf{97.0} \\
4 & 512 & 0.3 & 37.6 & 37.6 & 37.6 & 47.6 & 47.5 & 47.5 & 28.1 & 28.1 & 28.1 & \textbf{95.9} & \textbf{95.9} & \textbf{95.9} \\
\midrule
8 & 2 & 1.2 & \textbf{98.3} & \textbf{98.3} & \textbf{98.3} & 97.5 & 97.7 & 97.7 & 84.1 & 84.1 & 84.1 & 97.2 & 97.2 & 97.1 \\
8 & 4 & 1.2 & 88.4 & 88.4 & 88.4 & 89.4 & 89.4 & 89.4 & 90.6 & 90.6 & 90.6 & \textbf{96.7} & \textbf{96.7} & \textbf{96.7} \\
8 & 8 & 1.2 & 67.2 & 67.2 & 67.2 & 82.5 & 82.5 & 82.5 & 93.2 & 93.2 & 93.2 & \textbf{97.3} & \textbf{97.3} & \textbf{97.3} \\
8 & 16 & 1.0 & 53.3 & 53.3 & 53.3 & 77.3 & 77.2 & 77.2 & 94.2 & 94.2 & 94.2 & \textbf{97.1} & \textbf{97.1} & \textbf{97.1} \\
8 & 32 & 1.0 & 46.3 & 46.4 & 46.4 & 76.8 & 76.6 & 76.6 & 95.0 & 95.0 & 95.0 & \textbf{97.1} & \textbf{97.2} & \textbf{97.2} \\
8 & 64 & 0.8 & 39.8 & 39.8 & 39.8 & 72.4 & 72.8 & 72.8 & 93.1 & 93.1 & 93.1 & \textbf{96.9} & \textbf{97.0} & \textbf{97.0} \\
8 & 128 & 0.6 & 33.3 & 33.3 & 33.3 & 52.0 & 61.0 & 61.0 & 85.5 & 85.5 & 85.5 & \textbf{96.5} & \textbf{96.7} & \textbf{96.7} \\
8 & 256 & 0.5 & 27.6 & 27.6 & 27.6 & 47.5 & 47.5 & 47.5 & 51.2 & 51.2 & 51.2 & \textbf{95.9} & \textbf{95.9} & \textbf{95.9} \\
8 & 512 & 0.3 & 23.2 & 23.2 & 23.2 & 37.8 & 37.8 & 37.8 & 27.7 & 27.7 & 27.7 & \textbf{94.4} & \textbf{94.4} & \textbf{94.4} \\
\midrule
16 & 4 & 1.3 & 87.5 & 87.5 & 87.5 & 89.4 & 89.4 & 89.4 & 89.7 & 89.7 & 89.7 & \textbf{95.5} & \textbf{95.5} & \textbf{95.5} \\
16 & 8 & 1.4 & 60.9 & 60.9 & 60.9 & 78.8 & 79.2 & 79.2 & 91.2 & 91.2 & 91.2 & \textbf{95.6} & \textbf{95.6} & \textbf{95.6} \\
16 & 16 & 1.1 & 44.8 & 44.7 & 44.7 & 72.0 & 72.0 & 72.0 & 93.8 & 93.8 & 93.8 & \textbf{95.8} & \textbf{95.8} & \textbf{95.8} \\
16 & 32 & 1.0 & 35.9 & 35.9 & 35.9 & 72.0 & 71.3 & 71.3 & 94.2 & 94.2 & 94.2 & \textbf{95.5} & \textbf{95.6} & \textbf{95.6} \\
16 & 64 & 0.8 & 29.4 & 29.4 & 29.4 & 69.4 & 67.7 & 67.7 & 92.9 & 92.9 & 92.9 & \textbf{95.8} & \textbf{95.7} & \textbf{95.7} \\
16 & 128 & 0.6 & 23.2 & 23.2 & 23.2 & 58.5 & 55.5 & 55.5 & 85.4 & 85.4 & 85.4 & \textbf{94.6} & \textbf{94.7} & \textbf{94.7} \\
16 & 256 & 0.5 & 18.1 & 18.1 & 18.1 & 45.5 & 44.0 & 44.0 & 51.4 & 51.4 & 51.4 & \textbf{92.9} & \textbf{92.8} & \textbf{92.8} \\
16 & 512 & 0.3 & 14.3 & 14.3 & 14.3 & 29.1 & 29.1 & 29.1 & 30.4 & 30.4 & 30.4 & \textbf{90.1} & \textbf{90.1} & \textbf{90.2} \\
\bottomrule
\end{tabular}
\end{table}


\begin{algorithm}[htbp]
\caption{Evaluation protocol for one configuration (sampler, mixture variant, $K$, $d$, $T$, $b$).}
\label{alg:app-protocol}
\begin{lstlisting}[style=appcode]
# mu[0..K-1], w_mix : mode centres and weights, R99 : mode-core radius
# d_min : minimum centre separation, N : test seeds per mode, n_runs : runs
# F_star / F_theta : exact / learned field, G_theta : endpoint map of F_theta
# n, m : core and band anchors per mode, w : band width, T : pullback steps
# n_train, n_steps : training samples and steps, n_probe : probe seeds
# hall_max : retraining threshold, s_re : retraining step factor

def outcome(x, mu, R99):                     # k if x in C_k, -1 = hallucination
    D = cdist(x, mu)
    return where(min(D, axis=1) <= R99, argmin(D, axis=1), -1)

for run in range(n_runs):
    mu, w_mix = draw_centres(K, d, d_min), draw_weights(K)
    x_train   = sample_gmm(mu, w_mix, n_train)        # fixed training set
    F_theta   = train_field(x_train, n_steps)         # Heun/RK45/DPM++: DDIM's
    if hallucination_rate(G_theta, randn(n_probe, d)) > hall_max:
        F_theta = train_field(x_train, s_re * n_steps)   # once, same init

    z     = randn(N * K, d)                           # test seeds
    y_ref = outcome(G_theta(z), mu, R99)              # reference outcomes

    Z, y  = pullback(mu, R99, n, m, w, F_star, T)     # Alg. 1
    for clf in classifiers:                           # w-kNN: its own anchors
        y_hat = clf.fit(Z, y).predict(z)
        scores[run, clf] = metrics(y_hat, y_ref)      # accuracy, mode F1, hall. F1

return mean(scores, axis=0)
\end{lstlisting}
\end{algorithm}

The retraining rule is checked on $10^4$ probe seeds drawn independently of the test seeds, with
the process's own sampler (DDIM or flow matching); retraining restarts from the same initialisation
and training set with $1.7\times$ the steps and a correspondingly longer cosine schedule. Of the
420 DDIM and flow-matching networks (35 pairs $(K, d)$, two variants, three runs), three were
retrained, all at $K = 16$. One of them (DDIM, unweighted, $d = 512$, one run) remained above the
threshold, with a probe hallucination rate of $1.59\%$, and is kept and evaluated; every other
network ends with a probe rate of at most $1.48\%$. Heun, RK45 and DPM-Solver++(2M) reuse the DDIM
network unchanged, so the threshold does not bound their hallucination rate on the test seeds.

\begin{table}[htbp]
\centering
\small
\caption{\textbf{Experiment grid for the Gaussian mixtures.} Identifiers are those of
Algorithms~1 and~\ref{alg:app-protocol}.}
\label{tab:app-grid}
\begin{tabular}{lll}
\toprule
Setting & Identifier & Value \\
\midrule
Generative process           &                        & DDIM; flow matching \\
Solvers on the DDIM network  &                        & Heun; RK45; DPM-Solver++(2M) \\
Mixture weights              & \texttt{w\_mix}        & uniform; one $U[0.3, 0.8]$ profile, permuted per run \\
Number of modes              & \texttt{K}             & 2, 4, 8, 16 (not $K = 16$ at $d = 2$) \\
Dimension                    & \texttt{d}             & 2, 4, 8, 16, 32, 64, 128, 256, 512 \\
Minimum centre separation    & \texttt{d\_min}        & $3R_{99}$ \\
Pullback grid points         & \texttt{T}             & 250, 500, 750 \\
Core anchors per mode        & \texttt{n} $= b$       & 2{,}000, 5{,}000, 10{,}000, 15{,}000, 20{,}000 \\
Band anchors per mode        & \texttt{m}             & $b/2$ \\
Band width                   & \texttt{w}             & $2\sigma_0$ \\
Test seeds per mode          & \texttt{N}             & 50{,}000 (100{,}000 to 800{,}000 in total) \\
Runs                         & \texttt{n\_runs}       & 3 (seeds 0, 100, 200) \\
\bottomrule
\end{tabular}
\end{table}


\subsection{Diffusion Model on GMM}
\label{app:details-diffusion}

\paragraph{Samplers.}
We use two generative processes, and three further solvers for one of them. All samplers are
deterministic discretisations of Eq.~(1), so each seed has a single outcome. As in Sec.~3.1, the
seed is $\boldsymbol{z}$ at $\tau = 1$ and the sample is
$\boldsymbol{x} = \mathcal{G}_\theta(\boldsymbol{z})$ at $\tau = 0$.

\textit{DDIM.} The variance-preserving process of Sec.~B,
$\boldsymbol{X}_t = \sqrt{\bar\alpha_t}\, \boldsymbol{X}_0 + \sqrt{1 - \bar\alpha_t}\, \boldsymbol{\epsilon}$,
with $\beta$ linear in $[10^{-4}, 0.02]$ over 500 noise levels. The network is trained with the
standard $\boldsymbol{\epsilon}$-prediction objective at levels drawn uniformly from the 500 and
sampled with the deterministic DDIM update. This update is the Euler discretisation of the
probability-flow ODE, Eq.~(B.12), in the fixed-centre coordinates
$\boldsymbol{y} = \boldsymbol{x}/\sqrt{\bar\alpha_t}$ of Sec.~B with the added noise scale
$\eta_t = \sqrt{(1 - \bar\alpha_t)/\bar\alpha_t}$ as time, where it reads
$\mathrm{d}\boldsymbol{y}/\mathrm{d}\eta = \hat{\boldsymbol{\epsilon}}_\theta$: Eq.~(2)
reparameterised by $\eta$, with the network's noise prediction in place of the exact one. The
sampler takes 499 steps between the 500 levels, from $\bar\alpha = 6.4\times10^{-3}$
($\eta = 12.5$) at the seed to $\bar\alpha = 1 - 10^{-4}$ ($\eta = 0.01$), and returns the sample
there, without a final denoising step.

\textit{Flow matching.} The optimal-transport conditional path of Lipman et al.~(2023) with
$\sigma_{\min} = 10^{-4}$, written in model time,
$\boldsymbol{X}_\tau = (1 - \tau)\, \boldsymbol{x} + (\sigma_{\min} + (1 - \sigma_{\min})\tau)\, \boldsymbol{z}$.
The network regresses the conditional velocity
$\boldsymbol{x} - (1 - \sigma_{\min})\,\boldsymbol{z} = -\partial_\tau \boldsymbol{X}_\tau$ at
$\tau \sim U[0, 1]$, so it outputs $-\boldsymbol{F}_\theta$, and is sampled with 499 Euler steps on a
uniform grid of 500 points in $\tau$. Its step-index input is $\mathrm{round}(499(1 - \tau))$, which
is $0$ at the seed.

\textit{Heun, RK45 and DPM-Solver++(2M).} The trained DDIM network, integrated between the same
500 noise levels (499 steps) with a different solver. All three integrate
$\mathrm{d}\boldsymbol{y}/\mathrm{d}\eta = \hat{\boldsymbol{\epsilon}}_\theta$ from $\eta = 12.5$ to
$0.01$: Heun's method (a DDIM predictor with a trapezoidal corrector; second order, 2 network
evaluations per step), fixed-step Dormand--Prince RK45 (fifth order, 6 evaluations per step) and
DPM-Solver++(2M) (Lu et al., 2022; a second-order multistep method on the data prediction in
$\lambda = -\log\eta$ whose first step is a DDIM step; 1 evaluation per step). Heun and
DPM-Solver++(2M) query the network only at the 500 levels it was trained on, whereas the four
intermediate RK45 stages query it at fractional step indices, interpolated linearly in $\eta$
between neighbouring levels. These samplers are not retrained: they share the DDIM network, its
anchors and its calibration seeds, and differ from DDIM only in the reference outcomes, which
isolates the effect of discretising $\boldsymbol{F}_\theta$ on the pulled-back cores.

\paragraph{Network.}
Every learned field uses the residual MLP of Figure~\ref{fig:app-scorenet}: four pre-LayerNorm
residual blocks with LeakyReLU(0.2) activations and a sinusoidal embedding of the step index. The
hidden width is $n_h = 256$ for $d \le 64$ and $n_h = 1024$ for $d \ge 128$. The network is trained
with a mean-squared-error objective, Adam with a cosine learning-rate decay, and an exponential
moving average (EMA) of the weights; the EMA weights are the ones probed and sampled
(Table~\ref{tab:app-scorenet}). The retraining rule is described after
Algorithm~\ref{alg:app-protocol}.

\begin{figure}[htbp]
\centering
\tikzset{
  >=Stealth, font=\small,
  layer/.style={draw, rounded corners=1pt, minimum width=3.4mm, minimum height=11mm, inner sep=0pt},
  lin/.style={layer, fill=blue!25},
  nrm/.style={layer, fill=yellow!45},
  act/.style={layer, fill=orange!40, minimum height=7mm},
  emb/.style={layer, fill=green!30},
  blk/.style={layer, fill=violet!20, minimum width=5mm, minimum height=13mm},
  sum/.style={draw, circle, inner sep=0pt, minimum size=4.5mm},
  lname/.style={font=\scriptsize, align=center},
  dim/.style={font=\scriptsize, text=black!65},
}
\begin{tikzpicture}
  % main path
  \node (x) {$\boldsymbol{X}_\tau$};
  \node[lin, right=7mm of x, label={[lname]below:Linear}, label={[dim]above:$d \to n_h$}] (lin) {};
  \node[sum, right=28mm of lin] (plus) {$+$};
  \coordinate[right=11.5mm of plus] (bpos);
  \node[blk] at ($(bpos)+(3.6mm,3.6mm)$) (b1) {};
  \node[blk] at ($(bpos)+(2.4mm,2.4mm)$) (b2) {};
  \node[blk] at ($(bpos)+(1.2mm,1.2mm)$) (b3) {};
  \node[blk] at (bpos) (b4) {};
  \node[lname, below=1mm of b4] {Residual\\block};
  \node[dim, above=0.5mm of b1] {$\times 4$};
  \node[nrm, right=10mm of b4, label={[lname]below:Layer\\Norm}] (n) {};
  \node[act, right=4mm of n, label={[lname]below:$\phi$}] (a) {};
  \node[lin, right=4mm of a, label={[lname]below:Linear}, label={[dim]above:$n_h \to d$}] (o) {};
  \node[right=11mm of o, align=center] (out) {$\hat{\boldsymbol{\epsilon}}_\theta$ or $-\boldsymbol{F}_\theta$};
  \node[lname, below=0.5mm of out] {noise (DDIM)\\velocity (flow)};
  % time branch
  \node[below=21mm of x] (t) {$t$};
  \node[lname, below=0.5mm of t] {step\\index};
  \node[emb] at (lin |- t) (e) {};
  \node[lname, below=1mm of e] {Sinusoidal\\embedding};
  \node[dim, above=0.5mm of e] {128};
  \node[act] at (plus |- t) (ta) {};
  \node[lname, below=1mm of ta] {$\phi$};
  \node[lin] at ($(e)!0.5!(ta)$) (tl) {};
  \node[lname, below=1mm of tl] {Linear};
  \node[dim, above=0.5mm of tl] {$128 \to n_h$};
  % edges
  \draw[->] (x) -- (lin);
  \draw[->] (lin) -- node[dim, above] {$n_h$} (plus);
  \draw[->] (plus) -- (b4);
  \draw[->] (b4.east -| b1.east) -- (n);
  \draw[->] (n) -- (a);
  \draw[->] (a) -- (o);
  \draw[->] (o) -- (out);
  \draw[->] (t) -- (e);
  \draw[->] (e) -- (tl);
  \draw[->] (tl) -- (ta);
  \draw[->] (ta) -- node[dim, right] {$n_h$} (plus);
\end{tikzpicture}

\vspace{6mm}
\begin{tikzpicture}
  \node (in) {$\boldsymbol{\zeta}$};
  \node[nrm, right=8mm of in, label={[lname]below:Layer\\Norm}] (n) {};
  \node[act, right=5mm of n, label={[lname]below:$\phi$}] (a1) {};
  \node[lin, right=5mm of a1, label={[lname]below:Linear}, label={[dim]above:$n_h \to n_h$}] (l1) {};
  \node[act, right=5mm of l1, label={[lname]below:$\phi$}] (a2) {};
  \node[lin, right=5mm of a2, label={[lname]below:Linear}, label={[dim]above:$n_h \to n_h$}] (l2) {};
  \node[sum, right=6mm of l2] (plus) {$+$};
  \node[right=6mm of plus] (o) {$\boldsymbol{\zeta}'$};
  \coordinate (split) at ($(in.east)!0.5!(n.west)$);
  \draw[->] (in) -- (n);
  \draw[->] (n) -- (a1);
  \draw[->] (a1) -- (l1);
  \draw[->] (l1) -- (a2);
  \draw[->] (a2) -- (l2);
  \draw[->] (l2) -- (plus);
  \draw[->] (plus) -- (o);
  \draw[->, rounded corners=3pt] (split) -- ++(0,11mm) -| node[pos=0.25, lname, above] {skip connection} (plus.north);
  \fill (split) circle (0.6pt);
\end{tikzpicture}
\caption{\textbf{Network of the learned field $\boldsymbol{F}_\theta$.} Top: the full network with
hidden width $n_h$. The step index $t$ is embedded sinusoidally (dimension 128), mapped to
$\mathbb{R}^{n_h}$ by a linear layer and $\phi$, and added to the input projection; for DDIM it is
the level $t \in \{0, \dots, 499\}$ ($t = 499$ at the seed, fractional at the intermediate RK45
stages), and for flow matching it is $\mathrm{round}(499(1 - \tau))$ ($0$ at the seed). Bottom: one
pre-LayerNorm residual block acting on a hidden activation
$\boldsymbol{\zeta} \in \mathbb{R}^{n_h}$. Box colours give the layer type: linear (blue),
LayerNorm (yellow), activation $\phi = \mathrm{LeakyReLU}(0.2)$ (orange), embedding (green),
residual block (violet); grey text gives the dimensions.}
\label{fig:app-scorenet}
\end{figure}

\begin{table}[htbp]
\centering
\small
\caption{\textbf{Learned field $\boldsymbol{F}_\theta$: architecture and training.} Identifiers are
those of Figure~\ref{fig:app-scorenet} and Algorithm~\ref{alg:app-protocol}.}
\label{tab:app-scorenet}
\begin{tabular}{lll}
\toprule
Setting & Identifier & Value \\
\midrule
Hidden width                  & $n_h$                  & 256 for $d \le 64$; 1024 for $d \ge 128$ \\
Residual blocks               &                        & 4 \\
Time embedding                &                        & sinusoidal, dimension 128, of the step index \\
Activation                    & $\phi$                 & LeakyReLU(0.2) \\
Output                        &                        & $\hat{\boldsymbol{\epsilon}}_\theta$ (DDIM) or $-\boldsymbol{F}_\theta$ (flow matching) \\
\midrule
Training set                  & \texttt{n\_train}      & $10^5$ samples from $p_{\text{data}}$, fixed per run \\
Optimiser                     &                        & Adam, learning rate $10^{-3}$, cosine decay to $10^{-5}$ \\
Batch size                    &                        & 512 \\
Gradient clipping             &                        & norm 1 \\
Weight EMA                    &                        & decay $\min(0.999, (1 + s)/1001)$ at step $s$ \\
Training steps                & \texttt{n\_steps}      & $\lfloor 30{,}000\,(1 + d/16)(1 + K/16) \rfloor$ \\
Probe seeds                   & \texttt{n\_probe}      & $10^4$, independent of the test seeds \\
Retraining threshold          & \texttt{hall\_max}     & probe hallucination rate $> 1.5\%$ \\
Retraining                    & \texttt{s\_re}         & once, same initialisation, $1.7\times$ the steps \\
Sampling grid                 &                        & 500 points (499 steps) \\
\bottomrule
\end{tabular}
\end{table}


\subsection{Algorithm Details}
\label{app:algo-details}

We compare four classifiers that map a seed $\boldsymbol{z}$ to one of the $K + 1$ outcomes, given
in Algorithms~\ref{alg:app-knn}--\ref{alg:app-param-train} and configured as in
Table~\ref{tab:app-classifiers}. $k$NN, the quadratic classifier and the polar classifier are fit
on the core and band anchors of Sec.~\ref{app:anchor-class}. Weighted $k$NN (w-$k$NN) builds its
own anchors by pulling back points on 5 concentric spheres of radii
$0.3R_{99}, 0.6R_{99}, \dots, 1.5R_{99}$ around each $\boldsymbol{\mu}_k$, all labelled with the mode
$k$ and none with $-1$ (the outer two spheres lie outside the core $\mathcal{C}_k^\bullet$), and
detects hallucinations by the low confidence of its weighted vote. The quadratic and polar
classifiers differ only in their logits (Algorithms~\ref{alg:app-quadratic}
and~\ref{alg:app-polar}); both are trained with cross-entropy over the $K + 1$ classes, as
ensembles whose log-probabilities are summed (Algorithm~\ref{alg:app-param-train}).

\begin{algorithm}[htbp]
\caption{$k$NN: seed $\to$ outcome.}
\label{alg:app-knn}
\begin{lstlisting}[style=appcode]
# Z, y : pulled-back anchors, labels in {-1, 0..K-1};  z : test seeds
# n_nb : neighbours, eps : numerical floor
Y      = one_hot(y + 1, K + 1)                    # column 0 = hallucination
idx    = topk(cdist(z, Z), n_nb, largest=False)   # n_nb nearest anchors
lp     = log(mean(Y[idx], axis=1) + eps)          # log vote shares
margin = lp[:, 0] - max(lp[:, 1:], axis=1)        # hallucination margin
y_hat  = where(margin > 0, -1, argmax(lp[:, 1:], axis=1))
\end{lstlisting}
\end{algorithm}

\begin{algorithm}[htbp]
\caption{Weighted $k$NN (w-$k$NN) over mode-labelled anchors: seed $\to$ outcome.}
\label{alg:app-altered}
\begin{lstlisting}[style=appcode]
# mu, R99, n, F_star, T : as in Alg. 1;  outcome, G_star : as in the protocol
# n_sph : spheres per mode, r_max : outermost radius (units of R99)
# w_min : weight of the outermost sphere, n_nb : neighbours, temp : temperature
# n_cal : calibration seeds;  n_q, [q_lo, q_hi] : threshold candidates
P, y, om = [], [], []
for k in range(K):
    for j in range(1, n_sph + 1):
        P  += [mu[k] + (j / n_sph) * r_max * R99 * unit_vectors(n // n_sph, d)]
        y  += [k]
        om += [1 - (1 - w_min) * j / n_sph]       # linear decay to w_min
Za = transport(P, F_star, T)                      # integration loop of Alg. 1

def vote(z):
    idx   = topk(cdist(z, Za), n_nb, largest=False)
    s     = sum(om[idx] * one_hot(y[idx], K), 1) / sum(om[idx], 1)
    probs = softmax(s / temp)
    conf  = 1 - entropy(probs) / log(K)
    return probs, conf

# calibration seeds, independent of the test seeds, labelled by the exact map
z_cal = randn(n_cal, d)
y_cal = outcome(G_star(z_cal), mu, R99)
_, c  = vote(z_cal)
c_thr = argmax_over(quantiles(c, q_lo .. q_hi, n_q),
                    lambda th: hall_f1(c < th, y_cal == -1))

probs, conf = vote(z)
y_hat = where(conf < c_thr, -1, argmax(probs, axis=1))
\end{lstlisting}
\end{algorithm}

The calibration seeds are the same for every $T$ and $b$ of a run and are labelled by
$\mathcal{G}_\star$ on the same $T$-point grid as the pullback; \texttt{c\_thr} is recalibrated for
every run, $T$ and $b$.

\begin{algorithm}[htbp]
\caption{Quadratic classifier: logits of a seed.}
\label{alg:app-quadratic}
\begin{lstlisting}[style=appcode]
# z : seeds;  output : K + 1 logits, column 0 = hallucination
def forward(z):
    iu  = triu_indices(d, d)                      # pairs i <= j
    psi = concat([z, z[:, iu[0]] * z[:, iu[1]]])  # d + d(d+1)/2 features
    return Linear[d + d * (d + 1) // 2 -> K + 1](psi)
\end{lstlisting}
\end{algorithm}

\begin{algorithm}[htbp]
\caption{Polar classifier: logits of a seed.}
\label{alg:app-polar}
\begin{lstlisting}[style=appcode]
# z : seeds;  n_hid : hidden size, deg : polynomial degree;  output : K + 1 logits
# alpha0, alpha1, log_beta : learned scalars of the hallucination logit
def forward(z):
    r0    = sqrt(d)                               # typical seed radius
    r     = norm(z, axis=1, keepdims=True)
    r_off = r - r0                                # radial offset
    v     = z / r                                 # direction
    hid   = Linear[d + 1 -> n_hid]([v * sqrt(r0), r_off]) / sqrt(r0 + 1)

    phi = concat([hid ** i for i in range(1, deg + 1)])
    l   = Linear[deg * n_hid -> K](phi)           # mode logits

    t1, t2 = top2(l)
    l_h = alpha0 + alpha1 * r_off - exp(log_beta) * (t1 - t2)   # hallucination logit
    return concat([l_h, l])
\end{lstlisting}
\end{algorithm}

\begin{algorithm}[htbp]
\caption{Training and prediction of the quadratic and polar classifiers.}
\label{alg:app-param-train}
\begin{lstlisting}[style=appcode]
# Z, y : pulled-back anchors;  arch in {quadratic, polar};  z : test seeds
# n_ens : ensemble size, n_batch : batch size, lr, wd : learning rate and weight decay
# e_min, s_min : minimum epochs and minimum optimiser steps
nets = []
for e in range(n_ens):
    net    = Classifier(d, K, arch, seed=e)
    epochs = max(e_min, ceil(s_min / (len(Z) // n_batch)))
    train(net, Z, y + 1, loss=cross_entropy, epochs=epochs, batch=n_batch,
          opt=AdamW(lr=lr, weight_decay=wd), sched="one-cycle")
    nets += [net]

lp     = sum(log_softmax(net(z)) for net in nets)
margin = lp[:, 0] - max(lp[:, 1:], axis=1)        # hallucination margin
y_hat  = where(margin > 0, -1, argmax(lp[:, 1:], axis=1))
\end{lstlisting}
\end{algorithm}

Each epoch reshuffles the anchors and drops the last incomplete batch, so a network takes between
$3{,}000$ optimiser steps ($1{,}500$ epochs at $K = 2$, $b = 2{,}000$) and $7{,}020$ ($30$ epochs
at $K = 16$, $b = 20{,}000$); there is no validation split or early stopping, and the final weights
are used. The one-cycle schedule has the PyTorch defaults: the learning rate rises from
$\texttt{lr}/25$ to \texttt{lr} over the first $30\%$ of the steps and then decays to
$\texttt{lr}/(2.5\times10^{5})$, both along cosines. The linear layers have the PyTorch default
initialisation, the scalars \texttt{alpha0}, \texttt{alpha1} and \texttt{log\_beta} start at $0$,
and ensemble member $e \in \{0, 1, 2\}$ seeds its initialisation and batch order with $e$ in every
run.

The polar classifier uses the concentration of $\|\boldsymbol{z}\|_2$ around $\sqrt{d}$ for
$\boldsymbol{z} \sim \gamma$, so the outcome depends mainly on the direction
$\boldsymbol{v} = \boldsymbol{z}/\|\boldsymbol{z}\|_2$. Its hallucination logit grows as the two
largest mode logits approach each other, so hallucinations are predicted between two pulled-back
cores, in a band whose width depends on the radial offset. This matches the geometry of Sec.~3.3,
where the class boundaries arise from the contact of neighbouring cores.

\begin{table}[htbp]
\centering
\small
\setlength{\tabcolsep}{4pt}
\caption{\textbf{Classifier settings.} Identifiers are those of
Algorithms~\ref{alg:app-knn}--\ref{alg:app-param-train}.}
\label{tab:app-classifiers}
\begin{tabular}{llll}
\toprule
Classifier & Setting & Identifier & Value \\
\midrule
$k$NN          & neighbours                & \texttt{n\_nb}     & 10 \\
               & numerical floor           & \texttt{eps}       & $10^{-9}$ \\
\midrule
w-$k$NN        & neighbours                & \texttt{n\_nb}     & 10 \\
               & spheres per mode          & \texttt{n\_sph}    & 5, radii $0.3j\,R_{99}$, $j = 1, \dots, 5$ \\
               & anchors per sphere        & \texttt{n // n\_sph} & $b/5$ ($Kb$ in total; $1.5Kb$ for the others) \\
               & outermost radius          & \texttt{r\_max}    & 1.5 (units of $R_{99}$) \\
               & anchor weights            & \texttt{w\_min}    & 0.2; sphere $j$ has weight $1 - 0.8\,j/5$ \\
               & softmax temperature       & \texttt{temp}      & 0.1 \\
               & calibration seeds         & \texttt{n\_cal}    & 5{,}000, labelled by $\mathcal{G}_\star$ \\
               & threshold candidates      & \texttt{n\_q}, \texttt{q\_lo}, \texttt{q\_hi} & 100 quantiles in $[0.5\%, 50\%]$ \\
\midrule
Quadratic      & features                  &                    & $z_i$ and $z_i z_j$, $i \le j$ \\
\midrule
Polar          & hidden size, degree       & \texttt{n\_hid}, \texttt{deg} & 256, 3 \\
\midrule
Quadratic and  & optimiser                 & \texttt{lr}        & AdamW, one-cycle, peak $3 \times 10^{-3}$ \\
polar          & weight decay              & \texttt{wd}        & $10^{-4}$ \\
               & batch size                & \texttt{n\_batch}  & 2{,}048 \\
               & epochs, optimiser steps   & \texttt{e\_min}, \texttt{s\_min} & at least 30 and at least 3{,}000 \\
               & ensemble size             & \texttt{n\_ens}    & 3 \\
\bottomrule
\end{tabular}
\end{table}


\section{Additional Experiments}
\label{app:additional_experiments}


\subsection{Hallucination Mitigation Table (for all d, K)}
\label{app:hall-mitigation}

\input{tables/mitigation_table}

\paragraph{Mitigation Strength and Hallucination Reduction}
\jeffrey{A comparison between strength of mitigation and hallucination reduction in flow matching and DDIM}

\paragraph{Flow Matching vs DDIM}
\jeffrey{Flow matching should work under smaller mitigation strength compared to DDIM}

\paragraph{Trajectory Divergence from Original with time}
\jeffrey{Showing that thevend of t}


\subsection{Additional MNIST Results}
\label{app:mnist}

% Checked against ~/Desktop/iclr/experiment-2, run abc123 (code of commit 2be1abf, results in
% 24c85fc; decoded-sample grids from commit 746f4e1). Its current HEAD defaults differ (exact field
% integrated with Euler, sample grids ordered by clarity with at least 200 flagged seeds,
% Fashion-MNIST at d = 64, CIFAR-10 at d = 128 with T_train = 500). Fashion-MNIST and CIFAR-10
% were run in abc123 with the MNIST settings, but their results are not in the repo, so only MNIST
% is described here.

We test whether pullback predicts outcomes on real data, whose distribution is only approximately a
Gaussian mixture (Sec.~5.3). We train an autoencoder whose latent space is shaped to resemble the
mixtures of Sec.~\ref{app:details_gmm}, fit a Gaussian mixture to the latents, run a latent DDIM in
that space, and predict the outcome of its seeds as before.

\paragraph{Latent space.}
MNIST has $C = 10$ classes, and the latent dimension is $d = 32$. We train a convolutional
autoencoder with encoder $\mathrm{Enc}$ and decoder $\mathrm{Dec}$ for 300 epochs. The encoder sees
only images; the class labels $c_p$ enter only the loss, which imposes the geometry of
Sec.~\ref{app:details_gmm} on the latents $\boldsymbol{x}_p = \mathrm{Enc}(\mathrm{img}_p)$ of a
minibatch $p = 1, \dots, B$:
\begin{equation}
\begin{aligned}
    \mathcal{L} = \; & \frac{1}{B}\sum_{p=1}^{B} \|\mathrm{img}_p - \mathrm{Dec}(\boldsymbol{x}_p + \boldsymbol{\xi}_p)\|_2^2
    + \lambda_{\text{sep}}\, \mathbb{E}_{c_p \ne c_q} \Bigl[\max\!\Bigl(0, \tfrac{d_{\text{sep}} - \|\boldsymbol{x}_p - \boldsymbol{x}_q\|_2}{\sigma_0}\Bigr)^{2}\Bigr] \\
    & + \lambda_{\text{cov}}\, \mathbb{E}_{c} \Bigl\| \tfrac{\boldsymbol{\Sigma}_c}{\sigma_0^2} - \boldsymbol{I}_d \Bigr\|_F^2
    + \lambda_{\text{norm}}\, \mathbb{E}_{c} \Bigl[\max\!\Bigl(0, \tfrac{\|\boldsymbol{\mu}_c\|_2 - r_\mu}{\sigma_0}\Bigr)^{2}\Bigr],
\end{aligned}
\end{equation}
where the pixels are scaled to $[-1, 1]$,
$\boldsymbol{\xi}_p \sim \mathcal{N}(\boldsymbol{0}, (0.15\,\sigma_0)^2 \boldsymbol{I}_d)$, and
$\boldsymbol{\mu}_c$ and $\boldsymbol{\Sigma}_c$ are the minibatch mean and (unbiased) covariance of
the latents of class $c$ (Sec.~3.2). The separation term averages over the cross-class pairs that
violate the margin, the covariance term over the classes with at least $d + 2$ members in the
minibatch, and the norm term over the classes present. The constants are those of the Gaussian
mixtures: $\sigma_0 = 0.1$, $r_\mu = 2\sqrt{d/2}$ and $d_{\text{sep}} = 3R_{99}$ with
$R_{99} = \sigma_0\sqrt{\chi^2_d(0.99)}$ as in Sec.~\ref{app:details_gmm}; $\lambda_{\text{sep}} = 5$,
$\lambda_{\text{cov}} = 2$ and $\lambda_{\text{norm}} = 1$ are ramped in linearly over the first 5
epochs. The terms ask for latents of different classes to be at least $d_{\text{sep}}$ apart, for
each class to have covariance $\sigma_0^2\boldsymbol{I}_d$, and for the class means to lie within
radius $r_\mu$; the means themselves are not prescribed.

\paragraph{Fitted mixture and mode cores.}
The latent distribution is only approximately Gaussian, so we fit a mixture to the encoded training
set without labels: $k$-means with $K = 10$ clusters ($k$-means++ initialisation, 4 restarts) gives
the centres $\hat{\boldsymbol{\mu}}_k$, the cluster fractions give the weights $\hat w_k$, and the
mean over clusters of the per-dimension within-cluster variance gives $\hat\sigma_0^2$. Each cluster
has its own radius $\hat R_{99,k} = \hat\sigma_k \sqrt{\chi^2_d(0.99)}$, with $\hat\sigma_k^2$ its
per-dimension variance, and we use the largest, $\hat R_{99} = \max_k \hat R_{99,k}$ (at most $0.62$
over the three runs, against $R_{99} = 0.73$ for the target geometry). Labels are used only
afterwards, to name each cluster by its majority class. Every training latent lies inside its own
cluster's ball, while up to $2.3\%$ of the encoded test images lie outside every ball. We enlarge
the mode cores to $\hat{\mathcal{C}}_k^\bullet = B(\hat{\boldsymbol{\mu}}_k, \hat r)$ with
$\hat r = s_{\text{core}} \hat R_{99}$, $s_{\text{core}} = 1.25$, so that samples just outside a
fitted ball are not counted as hallucinations. The outcome of a seed is $k$ if its sample lies in
$\hat{\mathcal{C}}_k^\bullet$ and $-1$ otherwise; a hallucination is a latent sample outside every
enlarged core.

\paragraph{Diffusion model.}
The learned field $\boldsymbol{F}_\theta$ is a DDIM whose linear schedule
$\beta \in [10^{-4}, 0.02]$ is discretised into 100 noise levels (99 sampling steps, terminal
$\bar\alpha = 0.36$), so the seeds $\boldsymbol{z} \sim \gamma$ are not the terminal marginal of
the forward process. It uses the network of Figure~\ref{fig:app-scorenet} ($n_h = 256$), trained
on the encoded training images rather than on samples of the fitted mixture. We deliberately train
it for a tenth of the budget of Sec.~\ref{app:details-diffusion},
$3{,}000\,(1 + d/16)(1 + K/16) = 14{,}625$ steps, without retraining, so that it hallucinates often
enough to evaluate detection: on $4.3\%$ of the test seeds under $\mathcal{G}_\theta$, against
$0.004\%$ under the exact endpoint map of the fitted mixture.

\paragraph{Anchors, reference outcomes and evaluation.}
Anchors are sampled as in Sec.~\ref{app:anchor-class}, with $n = b$ points in each enlarged core
and $m = b/2$ in the band of width $w = 2\hat\sigma_0$ around it; the w-$k$NN spheres extend to
$r_{\max}\hat r$. Unlike the Gaussian mixtures, the exact field of the data is unknown, so the
anchors are pulled back by inverting the learned DDIM sampler: each reverse step starts from the
noise prediction at the current point and refines it with $n_{\text{fp}} = 2$ fixed-point
iterations (3 network evaluations per step), so that sampling from an anchor seed returns the anchor
up to the fixed-point residual. The w-$k$NN threshold is calibrated on $5{,}000$ seeds labelled by
$\mathcal{G}_\theta$. The test set consists of $N K = 100{,}000$ seeds ($N = 10{,}000$ per mode),
labelled by their outcome under $\mathcal{G}_\theta$. The classifiers and metrics are those of
Secs.~\ref{app:anchor-class} and~\ref{app:algo-details}. We add one baseline: each test seed is
mapped with the exact endpoint map $\hat{\mathcal{G}}_\star$ of the fitted mixture, integrated with
Heun-type DDIM steps on the linear schedule discretised into 200 levels (terminal
$\bar\alpha = 0.13$, against $0.36$ for $\mathcal{G}_\theta$), and its outcome is compared with that
under $\mathcal{G}_\theta$. This baseline reaches $85.6\%$ accuracy but a hallucination F1 of only
$0.002$, since $\hat{\mathcal{G}}_\star$ almost never hallucinates. As an independent check, a CNN
classifier trained on the ten classes reads the decoded samples, to verify that a sample in
$\hat{\mathcal{C}}_k^\bullet$ decodes to an image of the class of cluster $k$.
Algorithm~\ref{alg:app-mnist} summarises the procedure and Table~\ref{tab:app-mnist} gives the
values of its constants.

\paragraph{Decoded samples.}
To show what the predicted seed regions look like, we draw $100{,}000$ fresh seeds, fit each
classifier on the anchors at the largest budget, and predict the outcome of every seed. For each
predicted class we decode the samples $\mathrm{Dec}(\mathcal{G}_\theta(\boldsymbol{z}))$ of up to 10
seeds: for a cluster $k$, the first seeds (in seed order) predicted as $k$ whose samples lie within
$0.7\,\hat R_{99}$ of $\hat{\boldsymbol{\mu}}_k$; for the hallucination class, the first seeds
predicted as hallucinations whose samples lie beyond $\hat r$ from every centre. For each classifier
we also report the precision and recall of its hallucination predictions on these seeds, and the
distance of the flagged samples to their nearest centre in units of $\hat r$. A classifier that
flags fewer than 200 seeds would have its highest-margin seeds added until 200 are flagged; every
classifier flagged at least 637.

\pranav{A visual demonstration of samples which are hallucinations the definition of hallucinations and their corresponding generations.}

\begin{algorithm}[htbp]
\caption{Evaluation protocol in the MNIST latent space, for one anchor budget $b$.}
\label{alg:app-mnist}
\begin{lstlisting}[style=appcode]
# images, labels : MNIST, C = K = 10;  d : latent dimension;  N : test seeds per mode
# s_core : core enlargement;  n, m : core and band anchors per mode
# n_fp : fixed-point iterations;  n_steps : training steps of F_theta
# n_cal : calibration seeds;  n_runs : runs

for run in range(n_runs):
    Enc, Dec = train_autoencoder(images, labels, d)      # loss L above
    x_lat    = Enc(images)                               # latent training set
    mu_hat, w_hat, sigma0_hat = kmeans_mixture(x_lat, K) # no labels
    r_hat    = s_core * max(R99_per_cluster(x_lat, mu_hat))
    F_theta  = train_field(x_lat, n_steps)               # DDIM, 100 levels

    z      = randn(N * K, d)                             # test seeds
    y_ref  = outcome(G_theta(z), mu_hat, r_hat)          # reference outcomes
    y_base = outcome(G_star_hat(z), mu_hat, r_hat)       # exact map, 200 levels
    scores[run, "exact"] = metrics(y_base, y_ref)

    z_cal = randn(n_cal, d)                              # w-kNN calibration
    y_cal = outcome(G_theta(z_cal), mu_hat, r_hat)
    Z, y  = pullback(mu_hat, r_hat, n, m, 2 * sigma0_hat, invert(G_theta, n_fp))
    for clf in classifiers:                              # w-kNN: its own anchors
        y_hat = clf.fit(Z, y).predict(z)
        scores[run, clf] = metrics(y_hat, y_ref)

return mean(scores, axis=0)
\end{lstlisting}
\end{algorithm}

\begin{table}[htbp]
\centering
\small
\caption{\textbf{Settings in the MNIST latent space.} Identifiers are those of
Algorithm~\ref{alg:app-mnist}.}
\label{tab:app-mnist}
\begin{tabular}{lll}
\toprule
Setting & Identifier & Value \\
\midrule
Dataset                    &                         & MNIST ($C = 10$ classes) \\
Latent dimension           & \texttt{d}              & 32 \\
Autoencoder                &                         & convolutional, 300 epochs, AdamW, lr $8\times10^{-4}$, cosine decay \\
Autoencoder minibatch      & $B$                     & $\max(512,\, 2K(d + 2)) = 680$ \\
Fitted mixture             &                         & $k$-means, $K = 10$; $\hat\sigma_0$ from within-cluster variance \\
Core enlargement           & \texttt{s\_core}        & 1.25 \\
Learned sampler            & $\mathcal{G}_\theta$    & DDIM, 100 levels (99 steps, terminal $\bar\alpha = 0.36$) \\
Exact-field baseline       & $\hat{\mathcal{G}}_\star$ & Heun-type DDIM, 200 levels (terminal $\bar\alpha = 0.13$) \\
Network training           & \texttt{n\_steps}       & $3{,}000\,(1 + d/16)(1 + K/16) = 14{,}625$, no retraining \\
Pullback                   & \texttt{n\_fp}          & inversion of $\mathcal{G}_\theta$, 2 fixed-point iterations per step \\
Anchors per mode           & \texttt{n} $= b$        & 10{,}000, 50{,}000, 100{,}000 \\
Band anchors per mode      & \texttt{m}              & $b/2$, band width $2\hat\sigma_0$ \\
Test seeds per mode        & \texttt{N}              & 10{,}000 (100{,}000 in total) \\
Calibration seeds          & \texttt{n\_cal}         & 5{,}000, labelled by $\mathcal{G}_\theta$ \\
Runs                       & \texttt{n\_runs}        & 3 (seeds 0, 100, 200) \\
\bottomrule
\end{tabular}
\end{table}
